\section{Introduction}\label{sec:intro}
Cosmic-ray muons cross tens of metres of rock or concrete, and the fraction that survives depends on the material along their path. This has been used to measure tunnel overburden \citep{george1955}, search for chambers in pyramids \citep{alvarez1970,morishima2017} and image volcanoes \citep{tanaka2007}; see \citep{procureur2018,bonechi2020} for reviews. Transmission maps give line-of-sight information; combining viewpoints supports three-dimensional reconstruction. Previous demonstrations include portable-detector imaging of infrastructure-scale targets and a complete nuclear reactor \citep{chaiwongkhot2021,procureur2023}.

One detector position gives only an angular projection: it tells how much material lies along each direction, but not at what distance. Moving the detector adds parallax, although two views still constrain vertical extent poorly. Here we show that two positions of one portable tracker are enough for a three-dimensional voxel reconstruction of a building element: the beamed concrete ceiling of the cafeteria at Tel Aviv University. The voxel image is the main result. Ceiling-height diagnostics and a conditional beam-thickness fit support its interpretation and showcase the accompanying analysis. Ceiling height denotes the structure's vertical position; beam thickness is its extent along $z$. Tape-measured dimensions are independent comparisons, never fitting targets.

Cone-intersection localization and explicit comparison of reconstruction methods at the Temperino mine provide methodological context \citep{borselli2022}. We adapt geometric intersection and reprojection checks to angular count histograms; the single-view track-origin back-projection method used there is not implemented here.
